For $\nu=1/3,m=0$, the mean pointwise $n_1$ SEM changes from 0.0137 to 0.0084, and the $n_{\rm tot}$ SEM from 0.0246 to 0.0163. The five-band trace changes from $9.03\mathbin{\pm}0.2$ to $9.08\mathbin{\pm}0.071$.

For $\nu=2/3,m=1$, the mean pointwise $n_1$ SEM changes from 0.0197 to 0.0115, and the $n_{\rm tot}$ SEM from 0.0468 to 0.0249. The five-band trace changes from $18.1\mathbin{\pm}0.34$ to $17.8\mathbin{\pm}0.17$.

At $\nu=1/3$ the five new traces span 8.8502--9.0770.

At $\nu=2/3$ the five new traces span 17.6543--17.9703.

Across all ten states the mean pointwise SEM reduction is 16.8--62.4\% for $n_1$, and 14.3--72.4\% for $n_{\rm tot}$. This compares the complete old and new protocols (including different sample counts and eight versus 32 blocks), not an isolated test of importance sampling at equal cost. The number of nonlocal ratio evaluations per state increases from 528,384 to 1,320,960, a factor 2.5.

The maximum auxiliary orbital-normalization error is 0.00248; the largest change of a single band/momentum diagonal produced by this control is 0.00047. The uncorrected importance estimates remain in the released data for comparison. Small negative remote-band estimates are retained as evidence of finite-estimator uncertainty.
